\section{\ourloss{}: Fine-grained Regularized Medical Preference Optimization}
\label{section:method}

\subsection{Overview}
\label{sec:met:overview}

This chapter develops \ourloss{}, an alignment framework built directly against the four desiderata of Section~\ref{sec:ana:desiderata}. It comprises two halves that are trained jointly: a data construction procedure that produces preference pairs in which style is held constant and the clinically divergent spans are explicitly marked (D1, D2), and a training objective that combines a dual-constraint fine-grained ranking loss over clean and lesion-corrupted images (D1, D4) with a bidirectional token-wise KL regularizer (D3).

Before developing the components, it is worth stating precisely how \ourloss{} departs from RRPO, its closest antecedent. Three of RRPO's design choices do not survive the move to a clinical setting.

First, RRPO relies on a forward KL divergence term as its regularizer. This tends to preserve the coverage and diversity of $\pi_{\text{ref}}$, yet it often fails to facilitate the learning of new concepts that lie beyond the distribution of the reference model --- which is exactly what domain adaptation to medicine requires.

Second, RRPO imposes an overly rigid constraint by requiring the rejected response to strictly mirror a ground-truth template, which necessitates perfectly paired correct and incorrect phrases. In practical settings, such a constraint is often unrealistic, as a single phrase in the preferred response may correspond to multiple segments within the rejected response. Consequently, this forceful alignment can cause the model to shift its behaviour too abruptly when encountering new tasks. Rather than achieving seamless adaptation, this approach frequently leads to the forgetting of useful prior knowledge in favour of satisfying an off-policy alignment.

Third, the framework lacks explicit constraints regarding evidence grounding within the visual domain, a factor that is particularly critical for effective alignment in high-stakes environments such as the medical domain.

\subsection{On-Policy Fine-grained Preference Construction}
\label{sec:met:data}

The construction procedure takes an existing medical SFT dataset --- image $v$, question $q$, ground-truth answer $a$ --- and produces preference records in which $y^+$ and $y^-$ are stylistically matched and their clinically divergent spans are delimited by explicit \texttt{<mask>} markers. Figure~\ref{fig:pipeline_overview} gives the overall flow.

\begin{figure}[t]
    \centering
    \includegraphics[width=0.85\textwidth]{pipeline.pdf}
    \caption{An overview of the fine-grained on-policy visual-textual preference creation process.}
    \label{fig:pipeline_overview}
\end{figure}

\subsubsection{Sampling On-Policy Responses}

The procedure begins by obtaining a response from $\pi_\theta$ itself for the given $\{v, q\}$, using greedy decoding. This single decision is what makes the construction on-policy, and it is what distinguishes \ourloss{} from the reference-substitution approach analyzed in Section~\ref{sec:ana:reward-hacking}: both members of every resulting preference pair originate from the model's own output distribution, so the stylistic confound is eliminated by construction rather than corrected for after the fact.

Greedy decoding is chosen over sampling for a specific reason. The purpose of the sampled response is to expose the model's \emph{characteristic} behaviour --- what it will actually produce at inference time --- so that the correction targets an error the model genuinely makes. A high-temperature sample would exhibit errors drawn from the tail of the distribution, which the model would rarely commit in practice, and correcting those spends the training signal on failures that do not matter. The trade-off is coverage: greedy decoding yields one response per input, so the corpus is exactly as large as the underlying SFT dataset and no larger, whereas sampling could generate several pairs per input.

Because the response is the model's own, the corpus is specific to the backbone that generated it. Aligning a different model requires regenerating the textual side of the corpus. Section~\ref{sec:exp:corpus} quantifies the consequence: the two backbones yield corpora that differ in size and, more interestingly, in composition, since they make errors on different fractions of the data.

\subsubsection{Stage One: Localizing the Clinical Claim}

The sampled response is passed to a language model (\texttt{GPT-4o-mini}) together with the question and the ground-truth answer. The prompt, reproduced in full as Figure~\ref{fig:prompt_stage1} in Appendix~\ref{app:prompts}, asks for a single JSON object containing four fields:

\begin{itemize}
    \item \texttt{answer\_sentence} --- the minimal phrase or clause from the model's answer that states the final answer, in the model's exact wording.
    \item \texttt{final\_answer} --- the conclusive answer extracted in as few words as possible.
    \item \texttt{is\_correct} --- whether that answer is correct given the ground truth.
    \item \texttt{opposite\_answer} --- the answer sentence modified as little as possible into a medically plausible alternative.
\end{itemize}

The design of these fields is what enforces minimal editing. By requiring the extractor to quote the model's exact wording and to produce the counterfactual by modifying that quotation as little as possible, the resulting edit is confined to the clinical span; the surrounding sentence is carried over verbatim.

The \texttt{is\_correct} field then determines which branch of the construction applies:

\begin{itemize}
    \item If the output contains clinically designated errors, we substitute only the erroneous medical statements with ground-truth information to form $y^+$, while the original output serves as $y^-$.
    \item Conversely, if the initial response is already correct, we use it as $y^+$ and generate $y^-$ by altering the response with plausible medical counterfactuals --- the \texttt{opposite\_answer} field.
\end{itemize}

In both scenarios, the model's original linguistic style is preserved to allow on-policy alignment. The second branch matters more than it might appear: it means that questions the model already answers correctly still contribute a training signal, in the form of an explicit contrast against a plausible-but-wrong alternative, rather than being discarded.

\subsubsection{Stage Two: Localizing the Divergent Spans}

Given $y^+$ and $y^-$, a second call identifies exactly which phrases differ. The prompt presents the two sentences and asks the model to wrap only the differing phrases in \texttt{<mask>} and \texttt{</mask>} tags, with the instruction that the marked phrases be as short as possible.

These tags are the mechanism by which desideratum D1 is realized. They are added to the tokenizer as special tokens, and at training time the trainer converts them into a binary phrase mask over token positions --- every position strictly between a \texttt{<mask>} and its matching \texttt{</mask>} is marked, and the tag tokens themselves are then stripped from the input so the model never sees them. The ranking loss and the KL regularizer are both restricted to the marked positions. A single response may contain several such spans; the mask construction handles multiple pairs in one sequence.

\subsubsection{Adapting the Pipeline to Report Generation}

Radiology report generation does not fit the question-answer template, because there is no single answer span to localize --- a report contains many independent clinical assertions. The pipeline is therefore restructured for IU-Xray, using \texttt{Gemini 3 Flash Preview} in place of \texttt{GPT-4o-mini}.

Stage one becomes a constrained rewrite. The model's generated report (Report A) and the ground-truth report (Report B) are presented together, and the model is asked to use the facts in Report B to replace a few medical words in Report A. The prompt imposes three constraints that together preserve style while forcing a clinical change: \emph{structural locking} (the format of Report A must be preserved), \emph{token parsimony} (at most two or three words per sentence may be changed), and \emph{clinical significance} (synonym substitutions such as ``clear'' for ``normal'' are explicitly rejected). If every fact in Report B is already present in Report A, the model instead introduces a clinical error to create the rejected member of the pair. Stage two is then the same masked-difference step as in the VQA pipeline, applied to the two reports. Both report-side prompts are given as Figure~\ref{fig:prompt_stage2}.

\subsection{Lesion-Targeted Visual Corruption}
\label{sec:met:corruption}

The visual half of the construction produces, for each record, a corrupted counterpart $v'$ of the image. Section~\ref{sec:ana:grounding} argued that global perturbations are too blunt for the clinical case: what must be tested is not whether the model uses the image at all, but whether it uses the diagnostically decisive region. The corruption is therefore targeted at that region and leaves the surrounding anatomy intact.

The procedure has three steps. First, a short text prompt naming the clinical concept relevant to the query is extracted from each record by a language model --- for instance, the finding or anatomical structure that the question asks about. Second, this prompt is passed to MedSAM3-v1~\cite{liu2025medsam3}, a purely text-guided medical segmentation model, which returns segmentation masks for the named concept; where the model returns several detections, their union is taken as a single region. Third, zero-mean Gaussian noise is added per channel \emph{inside the mask only}, at a standard deviation of $45$ on the $0$--$255$ intensity scale, and the result is clipped back to the valid range.

Three properties of this choice of corruption are deliberate.

\emph{It destroys evidence without announcing itself.} Additive noise inside an anatomically plausible boundary leaves an image that still reads as a radiograph or CT slice. Cropping or blacking out a region produces an image that is trivially identifiable as degraded, which gives the model an alternative route to satisfying the objective: it can learn to detect the corruption and condition on that, rather than learning to depend on the lesion. Section~\ref{sec:res:abl-visual} shows both of those coarser strategies performing markedly worse, which is consistent with this concern.

\emph{It is local.} Everything outside the mask is bit-identical to the clean image. The paired comparison therefore isolates the contribution of the query-relevant region specifically, rather than the contribution of the image as a whole. This is the property that distinguishes the objective from the global perturbations used in prior multimodal preference work, as discussed in Section~\ref{sec:rel:multimodal}.

\emph{It is calibrated to remove the finding rather than the anatomy.} A standard deviation of $45$ on a $0$--$255$ scale is large enough to obliterate the texture differences that distinguish pathological from healthy tissue, but leaves gross structure discernible. The intent is to make the specific diagnostic claim unsupportable while keeping the image interpretable enough that a well-grounded model would produce a hedged rather than a nonsensical response.

Two fallbacks handle the cases where the pipeline cannot deliver a localized region. If no concept prompt could be extracted for a record, noise is applied to the whole image instead. If a prompt exists but MedSAM3 returns no mask, a lighter global perturbation ($\sigma = 12$) is applied. Each record carries a note recording which path was taken, so that fallback cases can be identified during analysis, and a flag indicating whether a usable corrupted image is available at all --- records without one simply do not contribute to the visual term of the loss. Section~\ref{sec:exp:corpus} reports how often each path is taken, and the figure is high enough to qualify the results: on SLAKE only $64.0\%$ of records receive genuinely lesion-targeted corruption.

Report generation is handled differently by design. A report-generation prompt asks for a description of the whole study rather than about one localized finding, so there is no single query-relevant region to segment, and holistic image-level noise is used throughout. The visual term therefore tests dependence on the image as a whole in that setting rather than on a specific finding, which is a weaker property than the one tested on the VQA datasets.

Section~\ref{sec:res:abl-visual} compares this lesion-targeted strategy against cropping, image-level noise, and full masking, and finds the localized variant substantially stronger on all three benchmarks.

\subsection{The Dual-Constraint Ranking Loss}
\label{sec:met:ranking}

Given the subtle visual nuances in the medical domain, our objective is to prevent the model from hallucinating correct-sounding responses when essential visual evidence is absent or corrupted. We achieve this by structuring the ranking loss around a dual constraint.

First, we define preference pairs $(q, v, y^+) \succ (q, v, y^-)$, which align $\pi_\theta$ toward medically accurate answers when supported by valid visual evidence $v$.

Second, we introduce pairs $(q, v', y^-) \succ (q, v', y^+)$, which compel the model to refrain from providing a definitive response when the underlying visual evidence is corrupted ($v'$). This ensures that the model prioritizes visual evidence over linguistic priors, effectively penalizing correct guesses made for the wrong reasons.

The reversal in the second pair is the crux of the design and deserves emphasis, because it is counterintuitive: the objective explicitly prefers the \emph{clinically wrong} answer under the corrupted image. It does so because the quantity being supervised is not correctness but \emph{warrant}. When the evidence supporting a specific claim has been destroyed, a model that still asserts that claim with undiminished confidence is demonstrating that the claim was never grounded in the image to begin with. Preferring $y^+$ under both $v$ and $v'$ would reward exactly that behaviour. Requiring the preference to \emph{flip} when the evidence is removed forces the model's confidence in the clinical span to be a function of the visual evidence rather than of a linguistic prior. Section~\ref{sec:res:abl-visual} tests this against the two natural alternatives --- preferring $(v, y^+)$ over $(v', y^+)$, and over $(v', y^-)$ --- and finds both weaker.

Note that in both cases, we provide a fine-grained reward only on the clinically divergent tokens between $y^+$ and $y^-$, that is, on the positions selected by the phrase mask of Section~\ref{sec:met:data}. The ranking loss is defined as:
\begin{equation}\label{eq:firempo-rank}
\begin{split}
\mathcal{L}^{\text{(Rank)}}_{\text{FiRe-MPO}}(\pi_{\theta};\pi_{\text{ref}}) = &
-\mathbb{E}
\biggl[
\log \sigma \left(\sum_{i} r_\theta(v,q,y^+_i) - \sum_{i} r_\theta(v,q,y^-_i)\right) \\
& + \gamma\log \sigma \left(\sum_{i} r_\theta(v',q,y^-_i) - \sum_{i} r_\theta(v',q,y^+_i)\right)
\biggr],
\end{split}
\end{equation}
where the index $i$ ranges over masked positions only and $\gamma$ weights the grounding term relative to the primary preference. Setting $\gamma = 0$ recovers a purely textual fine-grained objective, which is the \emph{w/o visual pref.} ablation of Section~\ref{sec:res:abl-visual}.

\subsection{The Bidirectional Token-wise KL Regularizer}
\label{sec:met:kl}

Our goal is to enable the policy model to preserve the rich, generalized prior knowledge of the base model while simultaneously allowing it to assimilate new, specialized domain knowledge that may lie outside the base model's original distribution.

As established in Section~\ref{sec:ana:kl}, the forward KL divergence $\mathbb{D}_{\text{KL}}(\pi_{\text{ref}} \parallel \pi_\theta)$, as employed in RRPO, is mean-seeking, forcing the policy model to maintain the broad coverage of the reference distribution. Although this successfully preserves the base model's diverse vision-language capabilities, it is overly restrictive: it heavily penalizes the policy for shifting its probability mass to learn the new concepts necessary for adapting to specific fields such as medicine. Conversely, the reverse KL divergence $\mathbb{D}_{\text{KL}}(\pi_\theta \parallel \pi_{\text{ref}})$ is inherently mode-seeking, allowing the model to learn new concepts through effective tail suppression. However, it often induces a severe lack of generation diversity.

To address these limitations, we replace the unidirectional constraint with a bidirectional, token-wise KL regularizer. This balances coverage preservation against tail suppression by simultaneously applying both forward and reverse KL constraints at the token level:
\begin{equation}\label{eq:firempo-kl}
\begin{split}
\mathcal{L}^{\text{(KL)}}_{\text{FiRe-MPO}}(\pi_\theta; \pi_{\text{ref}}) = &
- \mathbb{E} \biggl[ \sum_t \Big( \lambda \mathbb{D}_{\text{KL}} \big(\pi_{\text{ref}}(y_t \mid v,q,y_{<t}) \parallel \pi_\theta(y_t \mid v,q,y_{<t}) \big) \\
& + (1-\lambda) \mathbb{D}_{\text{KL}} \big(\pi_\theta(y_t \mid v,q,y_{<t}) \parallel \pi_{\text{ref}}(y_t \mid v,q,y_{<t}) \big) \Big) \biggr],
\end{split}
\end{equation}
where $\lambda$ is a hyperparameter that dictates the relative influence of each constraint. The endpoints recover the unidirectional cases: $\lambda = 1$ gives the forward-only regularizer of RRPO, and $\lambda = 0$ gives a purely reverse-KL objective. Both are ablated in Section~\ref{sec:res:abl-kl}, alongside removing the regularizer entirely.

The two terms are not redundant, and it is worth being precise about why a single coefficient can balance them. Because the divergences are evaluated between the same pair of distributions at the same token positions, they differ only in which distribution supplies the expectation. The forward term is dominated by positions where the reference is confident and the policy has moved away, and penalizes the policy for abandoning reference-supported continuations. The reverse term is dominated by positions where the policy has become confident in something the reference finds unlikely, and penalizes exactly the fluent-but-unsupported continuation that Section~\ref{sec:bg:kl} identifies as the characteristic medical failure. A given update can incur one penalty without the other, so the sum is not merely a rescaling of either term; $\lambda$ selects which kind of deviation the optimizer is more willing to tolerate.

Like the ranking loss, this sum is restricted to the masked phrase positions rather than applied over the full response. This is a deliberate pairing rather than an implementation convenience. The ranking term concentrates the preference signal on the clinical spans, which is precisely where the policy is being pushed hardest and therefore where uncontrolled drift is most likely; constraining the same positions places the regularizer where the pressure is. It also keeps the two terms commensurable, since both are sums over the same index set, which is what makes a single coefficient $\alpha$ meaningful in Equation~\ref{eq:firempo-total}.

\subsection{The Complete Objective}
\label{sec:met:objective}

The complete \ourloss{} objective unifies the ranking loss with the KL regularizer:
\begin{equation}\label{eq:firempo-total}
\mathcal{L}_{\text{FiRe-MPO}}(\pi_\theta; \pi_{\text{ref}}) = \mathcal{L}^{\text{(Rank)}}_{\text{FiRe-MPO}}(\pi_{\theta};\pi_{\text{ref}}) + \alpha \cdot \mathcal{L}_{\text{FiRe-MPO}}^{\text{(KL)}}(\pi_\theta; \pi_{\text{ref}}),
\end{equation}
where $\alpha$ is a loss coefficient controlling the relative influence of the fine-grained reward and the model divergence.

The objective therefore has four hyperparameters. $\beta$ is the standard DPO preference temperature inside $r_\theta$; $\gamma$ weights the visual grounding term; $\lambda$ sets the balance between the two KL directions; and $\alpha$ scales the regularizer against the ranking loss. Default values are $\beta = 0.1$, $\gamma = 0.1$, $\lambda = 0.5$, and $\alpha = 0.01$, with the full configuration given in Appendix~\ref{app:impl}. Appendix~\ref{app:pseudocode} gives the objective in executable form.

\subsection{Implementation}
\label{sec:met:implementation}

\subsubsection{Phrase Masking and Batch Construction}

The \texttt{<mask>} and \texttt{</mask>} delimiters are registered as special tokens on the tokenizer before training. During collation, each sequence is scanned for delimiter pairs and a binary mask is produced marking every position strictly between a \texttt{<mask>} and its matching \texttt{</mask>}; the delimiter tokens themselves are then removed from the input, together with their positions in the mask, so that the model is never conditioned on markup that will not appear at inference time. Sequences containing an unmatched \texttt{<mask>} are rejected rather than silently truncated.

Because the ranking loss in Equation~\ref{eq:firempo-rank} contains two preference terms over two different images, each training step requires forward passes over four sequence variants --- $(v, y^+)$, $(v, y^-)$, $(v', y^+)$, $(v', y^-)$ --- under both the policy and the frozen reference. The collator carries the corrupted image alongside the clean one, together with a per-example boolean indicating whether a corrupted image exists for that record; the visual term is masked out for examples where it does not, so that records whose lesion could not be segmented still contribute their textual preference.

\subsubsection{Parameter-Efficient Adaptation}

Following standard preference learning protocols, we use Low-Rank Adaptation (LoRA)~\cite{hu2022lora}. LoRA is applied exclusively to the LLM component of the LVLM, while the vision encoder and cross-modal projectors remain frozen to preserve the fundamental perceptual features learned during pre-training. This choice interacts with the grounding objective in a way worth making explicit: because the vision tower is frozen, the visual preference term cannot improve grounding by changing how the image is encoded. It can only change how the language model \emph{attends to and uses} the existing visual representation. The grounding improvements reported in Section~\ref{sec:res:grounding} are therefore attributable to the objective rather than to visual representation learning.

A practical benefit of LoRA in this setting is that the reference model does not need to be held separately in memory: reference log-probabilities are obtained by disabling the adapters on the same base weights.

\subsubsection{Computational Cost}

The dominant cost relative to standard DPO is the number of forward passes per optimization step. DPO requires two policy and two reference passes, over $(v, y^+)$ and $(v, y^-)$. \ourloss{} requires four of each, because of the additional pair over $v'$. The KL regularizer adds no forward passes --- it is computed from logits already materialized by the $(v, y^+)$ pass --- but it does require retaining full vocabulary-sized logits for both policy and reference at those positions in order to evaluate the divergence, which is the dominant memory term. Backward cost is unchanged in structure, since only the LoRA parameters receive gradients.

Table~\ref{tab:cost} reports the measured cost of a \ourloss{} training step on \texttt{Qwen3-VL-4B-Instruct} over the SLAKE corpus, on a single H200 with 143\,GB of memory. Two practical observations follow from these measurements.

\begin{table}[ht]
\centering
\small
\caption{Measured training cost of \ourloss{} on \texttt{Qwen3-VL-4B-Instruct} over the SLAKE preference corpus, at the batch size and LoRA configuration of Appendix~\ref{app:impl}. Steady-state step time excludes the first optimization step, which is roughly an order of magnitude slower due to kernel autotuning and allocator warmup.}
\label{tab:cost}
\begin{tabular}{lr}
\toprule
\textbf{Quantity} & \textbf{Measured} \\
\midrule
Forward passes per step (policy + reference) & $4 + 4$ \\
Steady-state step time & 4.3--4.6\,s \\
First-step time (warmup) & $\approx$74\,s \\
Steps per epoch (SLAKE, global batch 8) & 610 \\
Wall-clock per epoch & $\approx$44\,min \\
Peak GPU memory (gradient checkpointing on) & $\approx$139\,GB \\
Peak GPU memory (gradient checkpointing off) & exceeds 143\,GB \\
\bottomrule
\end{tabular}
\end{table}

First, gradient checkpointing is not optional at this configuration. Disabling it --- ordinarily an attractive trade on a large-memory device, since it removes the recomputation of activations during the backward pass --- exhausts the 143\,GB card. The reason is the combination of eight forward passes over sequences that each carry on the order of a thousand visual tokens, together with the vocabulary-sized logits the KL term requires. The peak of roughly 139\,GB with checkpointing enabled also means that two training runs cannot share one device, which constrains how a large sweep must be scheduled.

Second, the cost is dominated by the visual pair rather than by the regularizer. The extra $(v', y^{\pm})$ passes double the forward work relative to DPO, whereas the bidirectional KL term reuses logits that have already been computed. A practitioner constrained on compute should therefore reduce $\gamma$ or drop the visual term before touching $\alpha$; Section~\ref{sec:res:abl-kl} shows the regularizer to be the component whose removal is most damaging, and Section~\ref{sec:res:dynamics} shows it to be the cheapest.

The data construction pipeline additionally incurs a one-off cost, paid once per dataset rather than per training step: two language model calls per record for the textual pipeline, and one segmentation pass plus one prompt-extraction call per record for the visual pipeline. For SLAKE this amounts to roughly ten thousand language model calls and five thousand segmentation passes, after which the corpus is fixed and reusable across every subsequent training run.
